% !TEX program = lualatex
% Transparencias de Fundamentos de las Matemáticas II
% Clase 15: Números complejos

\documentclass[aspectratio=169,11pt]{beamer}

\usepackage{fontspec}
\usepackage[spanish,es-nodecimaldot]{babel}
\usepackage{amsmath,amssymb,mathtools}
\usepackage{booktabs}
\usepackage{csquotes}
\usepackage{xcolor}

\usetheme{Madrid}
\usecolortheme{default}
\setbeamertemplate{navigation symbols}{}
\setbeamertemplate{footline}[frame number]

\definecolor{FdMBlue}{RGB}{20,55,100}
\definecolor{FdMRed}{RGB}{130,30,30}
\definecolor{FdMGreen}{RGB}{25,100,60}

\setbeamercolor{structure}{fg=FdMBlue}
\setbeamercolor{title}{fg=white,bg=FdMBlue}
\setbeamercolor{frametitle}{fg=white,bg=FdMBlue}
\setbeamercolor{block title}{fg=white,bg=FdMBlue}
\setbeamercolor{block body}{bg=blue!3}
\setbeamercolor{alerted text}{fg=FdMRed}

\setmainfont{Latin Modern Roman}
\setsansfont{Latin Modern Sans}
\setmonofont{Latin Modern Mono}

\newcommand{\N}{\mathbb{N}}
\newcommand{\Z}{\mathbb{Z}}
\newcommand{\Q}{\mathbb{Q}}
\newcommand{\R}{\mathbb{R}}
\newcommand{\C}{\mathbb{C}}
\newcommand{\Pow}{\mathcal{P}}
\newcommand{\id}{\operatorname{id}}
\newcommand{\mcd}{\operatorname{mcd}}
\newcommand{\mcm}{\operatorname{mcm}}
\newcommand{\im}{\operatorname{Im}}

\title[FdM II -- Clase 15]{Números complejos}
\subtitle{Forma binómica, forma polar y ecuaciones algebraicas}
\author{Antonio Falcó Montesinos}
\institute{Universidad CEU Cardenal Herrera}
\date{Fundamentos de las Matemáticas II}

\begin{document}

\begin{frame}
  \titlepage
\end{frame}

\begin{frame}{Objetivos de la clase}
\begin{itemize}
  \item Motivar la unidad imaginaria.
  \item Operar con números complejos.
  \item Usar forma polar para potencias y raíces.
\end{itemize}
\end{frame}

\begin{frame}{Mapa de la clase}
\tableofcontents
\end{frame}

\section{Forma binómica}

\begin{frame}{Forma binómica}
\begin{block}{Unidad imaginaria}
\begin{itemize}
  \item Se introduce \(i\) con \(i^2=-1\).
  \item Esto no afirma que exista un real con cuadrado negativo.
  \item Amplía el sistema para resolver \(x^2+1=0\).
\end{itemize}
\end{block}
\end{frame}

\begin{frame}{Forma binómica}
\begin{block}{Complejos}
\begin{itemize}
  \item Un complejo se escribe \(z=a+bi\), con \(a,b\in\mathbb{R}\).
  \item \(a\) es la parte real y \(b\) la parte imaginaria.
  \item \(\mathbb{C}=\{a+bi:a,b\in\mathbb{R}\}\).
\end{itemize}
\end{block}
\end{frame}

\begin{frame}{Forma binómica}
\begin{block}{Operaciones}
\begin{itemize}
  \item \((a+bi)+(c+di)=(a+c)+(b+d)i\).
  \item \((a+bi)(c+di)=(ac-bd)+(ad+bc)i\).
  \item La fórmula del producto usa \(i^2=-1\).
\end{itemize}
\end{block}
\end{frame}

\section{Módulo y forma polar}

\begin{frame}{Módulo y forma polar}
\begin{block}{Conjugado y módulo}
\begin{itemize}
  \item \(\overline{a+bi}=a-bi\).
  \item \(|a+bi|=\sqrt{a^2+b^2}\).
  \item Se cumple \(z\overline z=|z|^2\).
\end{itemize}
\end{block}
\end{frame}

\begin{frame}{Módulo y forma polar}
\begin{block}{Inverso}
\begin{itemize}
  \item Si \(z\neq 0\), entonces \(z^{-1}=\overline z/|z|^2\).
  \item Esta fórmula permite dividir complejos.
  \item El denominador \(|z|^2\) es real positivo.
\end{itemize}
\end{block}
\end{frame}

\begin{frame}{Módulo y forma polar}
\begin{block}{Forma polar}
\begin{itemize}
  \item Si \(z\neq 0\), puede escribirse \(z=r(\cos\theta+i\sin\theta)\).
  \item \(r=|z|\).
  \item \(\theta\) es un argumento y no es único.
\end{itemize}
\end{block}
\end{frame}

\section{Potencias, raíces y ecuaciones}

\begin{frame}{Potencias, raíces y ecuaciones}
\begin{block}{De Moivre}
\begin{itemize}
  \item \((\cos\theta+i\sin\theta)^n=\cos(n\theta)+i\sin(n\theta)\).
  \item La forma polar simplifica potencias.
  \item También permite calcular raíces.
\end{itemize}
\end{block}
\end{frame}

\begin{frame}{Potencias, raíces y ecuaciones}
\begin{block}{Raíces}
\begin{itemize}
  \item Las raíces \(n\)-ésimas de \(w\neq 0\) se distribuyen angularmente.
  \item Aparecen \(n\) soluciones distintas.
  \item El caso \(w=-1\), \(n=2\), produce \(i\) y \(-i\).
\end{itemize}
\end{block}
\end{frame}

\begin{frame}{Potencias, raíces y ecuaciones}
\begin{block}{Teorema Fundamental del Álgebra}
\begin{itemize}
  \item Todo polinomio no constante con coeficientes complejos tiene al menos una raíz compleja.
  \item Por tanto, en \(\mathbb{C}\) los polinomios se descomponen completamente en factores lineales.
  \item Este resultado explica el papel central de \(\mathbb{C}\) en ecuaciones algebraicas.
\end{itemize}
\end{block}
\end{frame}

\section{Profundización}

\begin{frame}{Profundización}
\begin{block}{Error frecuente}
\begin{itemize}
  \item La parte imaginaria de \(a+bi\) es \(b\), no \(bi\).
  \item El argumento no es único.
  \item No se debe aplicar la fórmula cuadrática ignorando raíces complejas.
\end{itemize}
\end{block}
\end{frame}

\begin{frame}{Profundización}
\begin{block}{Ejemplo}
\begin{itemize}
  \item Resolver \(z^2+4=0\).
  \item Equivale a \(z^2=-4\).
  \item Las soluciones son \(z=2i\) y \(z=-2i\).
\end{itemize}
\end{block}
\end{frame}

\begin{frame}{Profundización}
\begin{block}{Cierre de FdM II}
\begin{itemize}
  \item Enteros: opuestos y división euclídea.
  \item Divisibilidad: \(\mcd\), Bézout y factorización.
  \item Modular: clases, inversos y TCR.
  \item Racionales, reales y complejos completan la cadena de sistemas numéricos.
\end{itemize}
\end{block}
\end{frame}


\section{Detalle y práctica}

\begin{frame}{Detalle y práctica}
\begin{block}{Raíces y argumento}
\begin{itemize}
  \item Si \(z=re^{i\theta}\), sus raíces \(n\)-ésimas tienen módulo \(\sqrt[n]{r}\).
  \item Los argumentos se reparten como \(\frac{\theta+2k\pi}{n}\), con \(k=0,\ldots,n-1\).
  \item Las raíces aparecen uniformemente distribuidas sobre una circunferencia.
\end{itemize}
\end{block}
\end{frame}

\begin{frame}{Detalle y práctica}
\begin{block}{Actividad breve}
\begin{itemize}
  \item Calcular las raíces cuadradas de \(-1\).
  \item Calcular las raíces cúbicas de \(1\).
  \item Interpretar geométricamente la conjugación \(z\mapsto \overline{z}\).
\end{itemize}
\end{block}
\end{frame}

\begin{frame}{Cierre}
\begin{itemize}
  \item Los complejos resuelven ecuaciones sin solución real.
  \item El conjugado y el módulo permiten dividir.
  \item La forma polar organiza potencias y raíces.
\end{itemize}
\end{frame}

\begin{frame}{Trabajo recomendado}
\begin{itemize}
  \item Releer las secciones correspondientes del manual.
  \item Identificar definiciones, símbolos nuevos y enunciados principales.
  \item Preparar dudas y ejemplos para el seminario asociado.
\end{itemize}
\end{frame}

\end{document}
